\documentclass[12pt,a4paper]{article}
\usepackage[utf8]{inputenc}
\usepackage[T5]{fontenc}
\usepackage{amsmath, amssymb}
\usepackage{graphicx}
\usepackage{ifthen}

\newboolean{showsolutions}
\setboolean{showsolutions}{true}

% --- ĐỊNH NGHĨA CÁC LỆNH ẢO CHO CÔNG CỤ LATEX2JSON CỦA WEBSITE ---
\newcommand{\doctitle}[1]{\begin{center}\Large\textbf{#1}\end{center}}
\newcommand{\docdesc}[1]{\textbf{Mô tả:} #1}
\newcommand{\docgrade}[1]{\textbf{Lớp:} #1}
\newcommand{\docstatus}[1]{\textbf{Trạng thái:} #1}
\newcommand{\doctype}[1]{\textbf{Loại:} #1}
\newcommand{\doctopics}[1]{\textbf{Chủ đề:} #1}

\newenvironment{textblock}{}

\newenvironment{lesson}[2]{
	\noindent\textbf{\Large #1}\\
	\textit{#2}
}{}

\newcommand{\image}[3]{
	\begin{center}
		\includegraphics[width=#1\textwidth]{#3} \\
		\textit{#2}
	\end{center}
}

\newenvironment{quiz}[1]{
	\noindent\textbf{\Large Kiểm tra: #1}
}{}

\newenvironment{mcq}[4]{
	\noindent\textbf{Câu hỏi:} #1 \\
	\ifthenelse{\boolean{showsolutions}}{
		\textit{(Đáp án đúng: #2, Điểm: #3) - Giải thích: #4}
	}{}
	\begin{itemize}
	}{
	\end{itemize}
}
\newcommand{\option}[1]{\item #1}

\newenvironment{truefalse}[3]{
	\noindent\textbf{Đúng/Sai:} #1 \\
	\ifthenelse{\boolean{showsolutions}}{
		\textit{(Điểm: #2) - Giải thích: #3}
	}{}
	\begin{itemize}
	}{
	\end{itemize}
}
\newcommand{\statement}[2]{\item [\textbf{#1}] #2}

\newcommand{\shortanswer}[4]{
    \noindent\textbf{Trả lời ngắn (Điểm: #3):}

    \noindent #1

	\ifthenelse{\boolean{showsolutions}}{
		\noindent\textit{Đáp án: #2 - Giải thích:}
		\noindent #4
	}{}
}

\newcommand{\essay}[3]{
    \noindent\textbf{Tự luận (Điểm: #2):}
    
    \noindent #1
    
	\ifthenelse{\boolean{showsolutions}}{
		\noindent\textbf{Gợi ý / Đáp án:}
		\noindent #3
	}{}
}

\begin{document}

\doctitle{ĐỀ 03 - KHẢO SÁT HÀM SỐ ÔN THI 2026 (CÂU 91 - 135)}
\docdesc{Tuyển tập câu hỏi trắc nghiệm Khảo sát hàm số từ các đề thi thử tốt nghiệp THPT năm học 2025 - 2026 có lời giải chi tiết.}
\docgrade{Lớp 12}
\docstatus{draft}
\doctype{normal}
\doctopics{ham-so-va-do-thi}

\begin{quiz}{Trắc nghiệm nhiều lựa chọn}

\begin{mcq}{(THPT Liên cấp đại học Hồng Đức 2026) Tiệm cận đứng của đồ thị hàm số $y = \frac{2x+4}{x-1}$ là}{0}{1}{Vì $\lim_{x \to 1^+} \frac{2x+4}{x-1} = +\infty$ nên $x = 1$ là tiệm cận đứng của đồ thị hàm số.}
	\option{$x = 1$.}
	\option{$x = -1$.}
	\option{$x = 2$.}
	\option{$x = -2$.}
\end{mcq}

\begin{mcq}{(THPT Liên cấp đại học Hồng Đức 2026) Hàm số nào dưới đây có đồ thị như đường cong trong hình bên?\image{0.6}{}{lt_1.png}}{2}{1}{Đây là hình ảnh đồ thị của hàm số bậc ba với hệ số $a > 0$, nên chọn đáp án C ($y = x^3 - 3x - 1$).}
	\option{$y = \frac{x^2-2x+3}{x-1}$.}
	\option{$y = \frac{x+1}{x-1}$.}
	\option{$y = x^3-3x-1$.}
	\option{$y = x^2+x-1$.}
\end{mcq}

\begin{mcq}{(THPT Nguyễn Khuyến - LTT HCM 2026) Hàm số nào dưới đây đơn điệu trên tập xác định của nó?}{0}{1}{Xét hàm số $y = \sqrt{x^3+x-5}$ có tập xác định $[a;+\infty)$ với $a \approx 1{,}516$.\\Ta có $y' = \frac{3x^2+1}{2\sqrt{x^3+x-5}} > 0, \forall x \in (a;+\infty)$.\\Vậy hàm số đồng biến trên tập xác định của nó.}
	\option{$y = \sqrt{x^3+x-5}$.}
	\option{$y = \frac{4x-2}{x-1}$.}
	\option{$y = x^4+x^3+x+11$.}
	\option{$y = \cot 2x$.}
\end{mcq}

\begin{mcq}{(THPT Lê Thánh Tông - HCM 2026) Hình vẽ bên là đồ thị của hàm số $y = \frac{ax+b}{cx+b}$. Đường tiệm cận đứng của đồ thị có phương trình là\image{0.6}{}{lt_2.png}}{0}{1}{Quan sát đồ thị, đường tiệm cận đứng của đồ thị là $x = 1$.}
	\option{$x = 1$.}
	\option{$x = 2$.}
	\option{$y = 1$.}
	\option{$y = 2$.}
\end{mcq}

\begin{mcq}{(THPT Lê Thánh Tông - HCM 2026) Tìm giá trị lớn nhất của hàm số $y = \frac{2x-1}{x+1}$ trên đoạn $[1;2]$.}{3}{1}{Ta có $y' = \frac{3}{(x+1)^2} > 0, \forall x \in [1;2]$ nên hàm số đồng biến trên đoạn $[1;2]$.\\Do đó $\max_{[1;2]} y = y(2) = 1$.}
	\option{$\max_{[1;2]} y = \frac{1}{2}$.}
	\option{$\max_{[1;2]} y = -\frac{1}{2}$.}
	\option{$\max_{[1;2]} y = -\frac{1}{3}$.}
	\option{$\max_{[1;2]} y = 1$.}
\end{mcq}

\begin{mcq}{(THPT Lê Thánh Tông - HCM 2026) Tiệm cận xiên của đồ thị hàm số $y = \frac{-x^2-2x+5}{x+2}$ là}{0}{1}{Ta có $y = \frac{-x^2-2x+5}{x+2} = -x + \frac{5}{x+2}$.\\Vì $\lim_{x \to \pm\infty} [y - (-x)] = \lim_{x \to \pm\infty} \frac{5}{x+2} = 0$ nên đường thẳng $y = -x$ là tiệm cận xiên của đồ thị hàm số đã cho.}
	\option{$y = -x$.}
	\option{$y = -x+1$.}
	\option{$y = x+2$.}
	\option{$x = -2$.}
\end{mcq}

\begin{mcq}{(THPT Lê Thánh Tông - HCM 2026) Cho hàm số có đồ thị như hình vẽ bên. Phát biểu nào sau đây sai?\image{0.6}{}{lt_3.png}}{2}{1}{Dựa vào đồ thị hàm số ta có các ý A, B, D đều đúng.\\Điểm cực đại của hàm số là $x = 2$ (hoặc giá trị cực đại là $y = 4$), phát biểu "Điểm cực đại của hàm số là 4" là sai (nhầm lẫn giữa điểm cực đại và giá trị cực đại).}
	\option{Hàm số đồng biến trên khoảng $(0;2)$.}
	\option{Hàm số nghịch biến trên khoảng $(2;+\infty)$.}
	\option{Điểm cực đại của hàm số là 4.}
	\option{Điểm cực tiểu của đồ thị hàm số là $(0;0)$.}
\end{mcq}

\begin{mcq}{(Cụm trường Nghệ An 2026) Cho hàm số $y = f(x)$ có bảng biến thiên như hình vẽ bên. Mệnh đề nào dưới đây đúng?\image{0.6}{}{lt_4.png}}{3}{1}{Dựa vào bảng biến thiên, giá trị cực đại của hàm số là $y_{\text{CĐ}} = 5$.}
	\option{$\min_{\mathbb{R}} y = 4$.}
	\option{$y_{\text{CT}} = 0$.}
	\option{$\max_{\mathbb{R}} y = 5$.}
	\option{$y_{\text{CĐ}} = 5$.}
\end{mcq}

\begin{mcq}{(Cụm trường Nghệ An 2026) Tiệm cận đứng của đồ thị hàm số $y = \frac{5}{x-1}$ là đường thẳng có phương trình}{3}{1}{Ta có $\lim_{x \to 1^+} y = +\infty$ nên $x = 1$ là đường tiệm cận đứng của đồ thị hàm số.}
	\option{$x = 5$.}
	\option{$y = 0$.}
	\option{$y = 1$.}
	\option{$x = 1$.}
\end{mcq}

\begin{mcq}{(Cụm trường Nghệ An 2026) Số giao điểm của đồ thị hai hàm số $y = x^2-3x-1$ và $y = x^3-1$ là}{0}{1}{Xét phương trình hoành độ giao điểm: $x^3 - 1 = x^2 - 3x - 1 \Leftrightarrow x^3 - x^2 + 3x = 0 \Leftrightarrow x(x^2 - x + 3) = 0 \Leftrightarrow x = 0$.\\Vậy đồ thị hai hàm số có đúng 1 giao điểm.}
	\option{$1$.}
	\option{$0$.}
	\option{$3$.}
	\option{$2$.}
\end{mcq}

\begin{mcq}{(Cụm trường Nghệ An 2026) Cho hàm số $y = f(x)$ xác định và liên tục trên khoảng $(-\infty;+\infty)$, có bảng biến thiên như hình sau:\image{0.6}{}{lt_5.png}Mệnh đề nào sau đây đúng?}{2}{1}{Dựa vào bảng biến thiên ta có hàm số đồng biến trên khoảng $(-\infty;-1)$ và $(1;+\infty)$. Vậy đáp án C đúng.}
	\option{Hàm số nghịch biến trên khoảng $(-\infty;-1)$.}
	\option{Hàm số nghịch biến trên khoảng $(-\infty;2)$.}
	\option{Hàm số đồng biến trên khoảng $(1;+\infty)$.}
	\option{Hàm số đồng biến trên khoảng $(-1;+\infty)$.}
\end{mcq}

\begin{mcq}{(Cụm trường Nghệ An 2026) Đường cong trong hình dưới là đồ thị của hàm số nào?\image{0.6}{}{lt_6.png}}{1}{1}{Đồ thị hàm số đã cho có dạng hàm số bậc ba $y = ax^3+bx^2+cx+d \Rightarrow y' = 3ax^2+2bx+c$.\\Đồ thị đi qua điểm $(0;1) \Rightarrow d = 1$.\\Đi qua điểm $(2;5) \Rightarrow 8a+4b+2c+d = 5$.\\Có điểm cực trị tại $x = 0 \Rightarrow y'(0) = c = 0$.\\Có điểm cực trị tại $x = 2 \Rightarrow y'(2) = 12a+4b+c = 0$.\\Giải hệ ta được $a = -1, b = 3, c = 0, d = 1 \Rightarrow y = -x^3+3x^2+1$.}
	\option{$y = -x^3+3x^2-4$.}
	\option{$y = -x^3+3x^2+1$.}
	\option{$y = x^3-3x^2+1$.}
	\option{$y = -x^3+2x^2-1$.}
\end{mcq}

\begin{mcq}{(THPT Bãi Cháy - Quảng Ninh 2026) Cho hàm số $y = f(x)$ có bảng biến thiên như hình vẽ dưới đây:\image{0.6}{}{lt_7.png}Đồ thị hàm số $y = f(x)$ có bao nhiêu đường tiệm cận?}{1}{1}{Dựa vào bảng biến thiên, ta có:\\- $\lim_{x \to 1^+} y = +\infty$ và $\lim_{x \to 1^-} y = -\infty$ nên đường thẳng $x = 1$ là đường tiệm cận đứng của đồ thị hàm số.\\- $\lim_{x \to -\infty} y = -1$ nên đường thẳng $y = -1$ là đường tiệm cận ngang của đồ thị hàm số.\\- $\lim_{x \to +\infty} y = 2$ nên đường thẳng $y = 2$ là đường tiệm cận ngang của đồ thị hàm số.\\Vậy đồ thị hàm số có 3 đường tiệm cận.}
	\option{$2$.}
	\option{$3$.}
	\option{$1$.}
	\option{$4$.}
\end{mcq}

\begin{mcq}{(THPT Bãi Cháy - Quảng Ninh 2026) Tích của giá trị nhỏ nhất và giá trị lớn nhất của hàm số $y = x+\frac{4}{x}$ trên $[1;3]$ bằng}{3}{1}{Trên $[1;3]$, $y' = 1 - \frac{4}{x^2} = 0 \Leftrightarrow x = 2$.\\Ta có $y(1) = 5, y(2) = 4, y(3) = \frac{13}{3} \Rightarrow \min_{[1;3]} y = 4, \max_{[1;3]} y = 5$.\\Tích của giá trị nhỏ nhất và giá trị lớn nhất của hàm số là $4\cdot 5 = 20$.}
	\option{$6$.}
	\option{$\frac{52}{3}$.}
	\option{$\frac{65}{3}$.}
	\option{$20$.}
\end{mcq}

\begin{mcq}{(THPT Bãi Cháy - Quảng Ninh 2026) Cho hàm số $y = f(x)$ có đồ thị như hình vẽ dưới đây:\image{0.6}{}{lt_8.png}Tiệm cận ngang của đồ thị hàm số là đường thẳng có phương trình}{3}{1}{Tiệm cận ngang của đồ thị hàm số là đường thẳng có phương trình $y = 1$.}
	\option{$x = -\frac{1}{2}$.}
	\option{$x = 1$.}
	\option{$y = -\frac{1}{2}$.}
	\option{$y = 1$.}
\end{mcq}

\begin{mcq}{(Liên trường Hà Nội 2026) Cho hàm số $y = f(x)$ liên tục trên $\mathbb{R}$ và có bảng xét dấu đạo hàm như sau:\image{0.6}{}{lt_9.png}Hàm số $y = f(x)$ đồng biến trên khoảng nào trong các khoảng sau đây?}{0}{1}{Nhìn bảng xét dấu ta thấy trong khoảng $(0;1)$ thì đạo hàm $y' > 0$, do đó hàm số đồng biến trên khoảng $(0;1)$.}
	\option{$(0;1)$.}
	\option{$(-3;0)$.}
	\option{$(2;+\infty)$.}
	\option{$(-\infty;-1)$.}
\end{mcq}

\begin{mcq}{(Liên trường Hà Nội 2026) Cho đồ thị hàm số $y = \frac{2x+1}{2-x}$ có đường tiệm cận ngang là đường thẳng có phương trình}{0}{1}{Đồ thị hàm số $y = \frac{2x+1}{2-x}$ có đường tiệm cận ngang là $y = \frac{2}{-1} = -2$.}
	\option{$y = -2$.}
	\option{$x = -2$.}
	\option{$y = 1$.}
	\option{$x = 2$.}
\end{mcq}

\begin{mcq}{(THPT Nguyễn Trung Thiên - Hà Tĩnh 2026) Tiệm cận đứng của đồ thị hàm số $y = \frac{2x-2}{x+1}$ là}{0}{1}{Hàm số $y = \frac{2x-2}{x+1}$ có tập xác định là $D = \mathbb{R}\setminus\{-1\}$. Ta có $\lim_{x \to (-1)^+} \frac{2x-2}{x+1} = -\infty$ (vì tử số tiến về $-4$ và mẫu số tiến về $0^+$). Do đó, đường thẳng $x = -1$ là tiệm cận đứng của đồ thị hàm số.}
	\option{$x = -1$.}
	\option{$x = 1$.}
	\option{$x = -2$.}
	\option{$x = 2$.}
\end{mcq}

\begin{mcq}{(THPT Lê Thánh Tông - HCM 2026) Tìm hệ số $b, c$ để hàm số $y = \frac{2}{cx+b}$ có đồ thị như hình vẽ sau:\image{0.6}{}{lt_10.png}}{3}{1}{Đồ thị hàm số có tiệm cận đứng là $x = -\frac{b}{c} = 2 \Rightarrow b = -2c$.\\Đồ thị đi qua điểm $(0;-1) \Rightarrow \frac{2}{b} = -1 \Rightarrow b = -2 \Rightarrow c = 1$.\\Vậy $b = -2, c = 1$.}
	\option{$\begin{cases} b=2 \\ c=-1 \end{cases}$.}
	\option{$\begin{cases} b=1 \\ c=-1 \end{cases}$.}
	\option{$\begin{cases} b=2 \\ c=1 \end{cases}$.}
	\option{$\begin{cases} b=-2 \\ c=1 \end{cases}$.}
\end{mcq}

\begin{mcq}{(THPT Lê Thánh Tông - HCM 2026) Cho hàm số $y = f(x)$ có đồ thị như hình vẽ dưới đây\image{0.6}{}{lt_11.png}Hàm số đã cho đồng biến trên khoảng nào trong các khoảng sau đây?}{2}{1}{Từ đồ thị hàm số, ta thấy nhánh đồ thị đi lên từ trái sang phải trên $(2;+\infty)$, do đó hàm số đồng biến trên $(2;+\infty)$.}
	\option{$(-\infty;1)$.}
	\option{$(1;2)$.}
	\option{$(2;+\infty)$.}
	\option{$(-1;1)$.}
\end{mcq}

\begin{mcq}{(THPT Nguyễn Trãi - Hà Nội 2026) Cho hàm số $y = f(x)$ có bảng xét dấu đạo hàm như sau\image{0.6}{}{lt_12.png}Hàm số đã cho nghịch biến trên khoảng nào dưới đây?}{1}{1}{Hàm số đã cho nghịch biến trên khoảng $(3;4)$ vì $y' < 0$ với mọi $x \in (3;4)$.}
	\option{$(-\infty;2)$.}
	\option{$(3;4)$.}
	\option{$(4;+\infty)$.}
	\option{$(1;3)$.}
\end{mcq}

\begin{mcq}{(THPT Nguyễn Trãi - Hà Nội 2026) Cho hàm số $y = f(x)$ liên tục trên $\mathbb{R}$ và có bảng xét dấu của đạo hàm như sau:\image{0.6}{}{lt_13.png}Số điểm cực trị của hàm số đã cho là}{1}{1}{Dựa vào bảng xét dấu, đạo hàm đổi dấu 4 lần khi qua các điểm $x = -3, x = -2, x = 3, x = 5$. Do đó số điểm cực trị của hàm số đã cho là 4.}
	\option{$3$.}
	\option{$4$.}
	\option{$2$.}
	\option{$5$.}
\end{mcq}

\begin{mcq}{(THPT Nguyễn Trãi - Hà Nội 2026) Hình vẽ sau đây là đồ thị của hàm số nào trong các hàm số dưới đây?\image{0.6}{}{lt_14.png}}{1}{1}{Dựa vào hình vẽ suy ra hệ số $a > 0$ (Loại C).\\Đồ thị của hàm số có 2 điểm cực trị nên phương trình $y' = 0$ có 2 nghiệm phân biệt (Loại A).\\Đồ thị của hàm số giao với trục $Ox$ là $y = 0 \Leftrightarrow x^3-2x^2+1=0 \Leftrightarrow x = 1, x \approx -0{,}6, x \approx 1{,}6$ (Loại D).\\Vậy chọn B ($y = x^3 - 2x + 1$).}
	\option{$y = x^3+2x+1$.}
	\option{$y = x^3-2x+1$.}
	\option{$y = -x^3+2x+1$.}
	\option{$y = x^3-2x^2+1$.}
\end{mcq}

\begin{mcq}{(THPT Nguyễn Trãi - Hà Nội 2026) Cho hàm số $y = f(x)$ có đồ thị như hình vẽ bên dưới\image{0.6}{}{lt_15.png}Hàm số đã cho đồng biến trên khoảng nào dưới đây?}{2}{1}{Dựa vào đồ thị ta thấy, đồ thị hàm số đi lên từ trái sang phải trong khoảng $(2;+\infty)$ nên hàm số đã cho đồng biến trên khoảng $(2;+\infty)$.}
	\option{$(0;2)$.}
	\option{$(0;+\infty)$.}
	\option{$(2;+\infty)$.}
	\option{$(-\infty;2)$.}
\end{mcq}

\begin{mcq}{(THPT Nguyễn Trãi - Hà Nội 2026) Đồ thị hàm số $y = \frac{x+3}{x-1}$ có đường tiệm cận đứng là}{2}{1}{Ta có: Tập xác định: $D = \mathbb{R}\setminus\{1\}$.\\Vì $\lim_{x \to 1^+} \frac{x+3}{x-1} = +\infty$ nên đường thẳng $x = 1$ là đường tiệm cận đứng của đồ thị hàm số đã cho.}
	\option{$y = 1$.}
	\option{$x = -3$.}
	\option{$x = 1$.}
	\option{$x = -1$.}
\end{mcq}

\begin{mcq}{(THPT Nguyễn Trãi - Hà Nội 2026) Cho hàm số $y = f(x)$ có đồ thị trên đoạn $[-2;2]$ như hình vẽ\image{0.6}{}{lt_16.png}Gọi giá trị lớn nhất và giá trị nhỏ nhất của hàm số trên đoạn $[-2;2]$ lần lượt là $M$ và $m$. Khi đó $M-m$ bằng}{0}{1}{Ta có: $M = 1, m = -4$. Vậy $M - m = 1 - (-4) = 5$.}
	\option{$5$.}
	\option{$-4$.}
	\option{$0$.}
	\option{$3$.}
\end{mcq}

\begin{mcq}{(THPT Trần Nhân Tông - Hà Nội 2026) Cho hàm số $y = \frac{ax+b}{cx+d}$ ($ac \ne 0; ad-bc \ne 0$) có đồ thị như hình vẽ dưới đây. Trong các hệ số $a, b, c, d$ có bao nhiêu số dương?\image{0.6}{}{lt_17.png}}{0}{1}{Tiệm cận đứng là: $x = 1$ nên $-\frac{d}{c} = 1 \Rightarrow d = -c$.\\Tiệm cận ngang là: $y = -1$ nên $\frac{a}{c} = -1 \Rightarrow a = -c$.\\Đồ thị cắt trục tung tại điểm có tung độ bằng $-2$ nên $\frac{b}{d} = -2 \Rightarrow b = -2d = 2c$.\\Ta xét 2 trường hợp dấu của $c$:\\- TH1: Nếu $c > 0$ thì $b > 0; a < 0; d < 0$ suy ra có 2 số dương.\\- TH2: Nếu $c < 0$ thì $b < 0; a > 0; d > 0$ suy ra có 2 số dương.\\Vậy trong mọi trường hợp đều có 2 số dương.}
	\option{$2$.}
	\option{$1$.}
	\option{$3$.}
	\option{$0$.}
\end{mcq}

\begin{mcq}{(THPT Trần Phú - Hà Nội 2026) Cho hàm số $y = f(x)$ có bảng biến thiên như hình vẽ. Hàm số đã cho có bao nhiêu điểm cực trị.\image{0.6}{}{lt_18.png}}{0}{1}{Dựa vào bảng biến thiên trên thì hàm số có hai điểm cực trị ($x = -1$ và $x = 1$).}
	\option{$2$.}
	\option{$3$.}
	\option{$1$.}
	\option{$0$.}
\end{mcq}

\begin{mcq}{(THPT Trần Phú - Hà Nội 2026) Cho hàm số $y = f(x)$ liên tục trên $[-3;2]$ và có bảng biến thiên như hình dưới đây. Gọi $M$ và $m$ lần lượt là giá trị lớn nhất và giá trị nhỏ nhất của hàm số $y = f(x)$ trên $[-1;2]$. Giá trị của $M+m$ bằng bao nhiêu\image{0.6}{}{lt_19.png}}{1}{1}{Dựa vào bảng biến thiên của hàm số, trên đoạn $[-1;2]$ ta thấy hàm số đạt GTLN tại $x = -1$ là $M = 3$, hàm số đạt GTNN tại $x = 0$ là $m = 0$. Suy ra $M + m = 3 + 0 = 3$.}
	\option{$2$.}
	\option{$3$.}
	\option{$4$.}
	\option{$1$.}
\end{mcq}

\begin{mcq}{(Cụm Hải Phòng 2026) Cho hàm số $y = f(x)$ có bảng biến thiên như sau.\image{0.6}{}{lt_20.png}Hàm số $y = f(x)$ đồng biến trên khoảng nào sau đây?}{2}{1}{Từ bảng biến thiên, hàm số $y = f(x)$ đồng biến trên các khoảng $(1;2)$ và $(2;3)$.}
	\option{$(2;4)$.}
	\option{$(3;+\infty)$.}
	\option{$(1;2)$.}
	\option{$(0;1)$.}
\end{mcq}

\begin{mcq}{(Cụm Hải Phòng 2026) Cho hàm số $f(x)$ liên tục trên đoạn $[0;3]$ và có đồ thị như hình vẽ bên. Gọi $M$ và $m$ lần lượt là giá trị lớn nhất và nhỏ nhất của hàm số đã cho trên đoạn $[0;3]$. Giá trị của $M+m$ bằng\image{0.6}{}{lt_21.png}}{3}{1}{Dựa vào đồ thị, ta có: $M = \max_{[0;3]} f(x) = f(3) = 3$ và $m = \min_{[0;3]} f(x) = f(2) = -2 \Rightarrow M + m = 3 - 2 = 1$.}
	\option{$3$.}
	\option{$5$.}
	\option{$2$.}
	\option{$1$.}
\end{mcq}

\begin{mcq}{(THPT Lê Quý Đôn - Đống Đa 2026) Giá trị lớn nhất của hàm số $y = \frac{2x-5}{x-1}$ trên đoạn $[2;5]$ là}{2}{1}{Ta có $y' = \frac{3}{(x-1)^2} > 0, \forall x \in [2;5]$ nên hàm số đồng biến trên đoạn $[2;5] \Rightarrow \max_{[2;5]} y = y(5) = \frac{5}{4}$.}
	\option{$-1$.}
	\option{$\frac{1}{2}$.}
	\option{$\frac{5}{4}$.}
	\option{$1$.}
\end{mcq}

\begin{mcq}{(THPT Lê Quý Đôn - Đống Đa 2026) Giá trị nhỏ nhất của hàm số: $y = 4\sin 2x - 3$ là}{2}{1}{Ta có $-1 \le \sin 2x \le 1 \Rightarrow -4 \le 4\sin 2x \le 4 \Rightarrow -7 \le 4\sin 2x - 3 \le 1$.\\Vậy giá trị nhỏ nhất của hàm số $y = 4\sin 2x - 3$ là $-7$.}
	\option{$1$.}
	\option{$-11$.}
	\option{$-7$.}
	\option{$5$.}
\end{mcq}

\begin{mcq}{(THPT Lê Quý Đôn - Đống Đa 2026) Cho hàm số $y = f(x)$ liên tục trên $\mathbb{R}$ và có bảng biến thiên như sau\image{0.6}{}{lt_22.png}Hàm số đạt cực đại tại điểm}{3}{1}{Dựa vào bảng biến thiên, ta thấy hàm số $y = f(x)$ đạt cực đại tại điểm $x = 0$.}
	\option{$x = 1$.}
	\option{$x = -3$.}
	\option{$x = -2$.}
	\option{$x = 0$.}
\end{mcq}

\begin{mcq}{(THPT Lê Quý Đôn - Đống Đa 2026) Cho đồ thị hàm số như hình vẽ bên dưới. Đường tiệm cận ngang của đồ thị hàm số là\image{0.6}{}{lt_23.png}}{3}{1}{Đường tiệm cận ngang của đồ thị hàm số là $y = 1$.}
	\option{$x = 2$.}
	\option{$y = 2$.}
	\option{$x = 1$.}
	\option{$y = 1$.}
\end{mcq}

\begin{mcq}{(Sở Lạng Sơn 2026) Cho hàm số $y = f(x)$ có đồ thị như bên dưới\image{0.6}{}{lt_24.png}Hàm số đã cho nghịch biến trên khoảng nào trong các khoảng sau đây?}{2}{1}{Từ đồ thị ta thấy hàm số nghịch biến trên khoảng $(-1;0)$.}
	\option{$(0;2)$.}
	\option{$(-1;1)$.}
	\option{$(-1;0)$.}
	\option{$(1;2)$.}
\end{mcq}

\begin{mcq}{(Sở Thái Nguyên 2026) Tiệm cận ngang của đồ thị hàm số $y = \frac{2x-5}{x+2}$ là đường thẳng có phương trình}{0}{1}{Ta có $\lim_{x \to \pm\infty} y = \lim_{x \to \pm\infty} \frac{2x-5}{x+2} = 2$ suy ra đường thẳng $y = 2$ là tiệm cận ngang của đồ thị hàm số.}
	\option{$y = 2$.}
	\option{$y = -2$.}
	\option{$y = \frac{5}{2}$.}
	\option{$y = \frac{2}{5}$.}
\end{mcq}

\begin{mcq}{(Sở Thái Nguyên 2026) Cho hàm số $f(x)$ có bảng biến thiên như hình vẽ. Giá trị cực tiểu của hàm số đã cho là\image{0.6}{}{lt_25.png}}{3}{1}{Dựa vào bảng biến thiên ta thấy dấu của $f'(x)$ đổi dấu từ âm sang dương khi qua $x = 0$ nên giá trị cực tiểu của hàm số là $y_{\text{CT}} = f(0) = 1$.}
	\option{$y = 2$.}
	\option{$y = -1$.}
	\option{$y = 0$.}
	\option{$y = 1$.}
\end{mcq}

\begin{mcq}{(Sở Thái Nguyên 2026) Cho hàm số đa thức bậc bốn $y = f(x)$ có đồ thị như hình vẽ. Hàm số đã cho đồng biến trên khoảng nào dưới đây?\image{0.6}{}{lt_26.png}}{1}{1}{Quan sát đồ thị ta thấy hàm số đồng biến trên các khoảng $(-\infty;-1)$ và $(0;1)$.}
	\option{$(-1;+\infty)$.}
	\option{$(-\infty;-1)$.}
	\option{$(-1;0)$.}
	\option{$(-\infty;1)$.}
\end{mcq}

\begin{mcq}{(Sở Thái Nguyên 2026) Đồ thị của hàm số nào dưới đây có dạng như đường cong trong hình vẽ?\image{0.6}{}{lt_27.png}}{3}{1}{Dựa vào hình dạng đồ thị ta có: đồ thị là của hàm số bậc 3: $y = ax^3+bx^2+cx+d$ và có hệ số $a > 0$ nên chọn đáp án D ($y = x^3 - 3x + 1$).}
	\option{$y = \frac{2x-1}{x+2}$.}
	\option{$y = \frac{x^2-1}{2x-1}$.}
	\option{$y = -x^3+3x+1$.}
	\option{$y = x^3-3x+1$.}
\end{mcq}

\begin{mcq}{(Cụm trường Thanh Hóa 2026) Đường tiệm cận ngang đồ thị hàm số $y = \frac{3x-1}{x-2}$ là}{2}{1}{Ta có $\lim_{x \to \pm\infty} y = \lim_{x \to \pm\infty} \frac{3x-1}{x-2} = 3$ suy ra đường thẳng $y = 3$ là tiệm cận ngang của đồ thị hàm số.}
	\option{$y = 2$.}
	\option{$x = 3$.}
	\option{$y = 3$.}
	\option{$x = 2$.}
\end{mcq}

\begin{mcq}{(Cụm trường Thanh Hóa 2026) Cho hàm số $y = f(x)$ có đạo hàm $f'(x) = (x-1)(x-2)^2(x-3)^3$. Số điểm cực trị của hàm số $y = f(x)$ là}{0}{1}{\image{0.6}{}{lt_28.png}}
	\option{$2$.}
	\option{$3$.}
	\option{$1$.}
	\option{$0$.}
\end{mcq}

\begin{mcq}{(THPT Lê Thánh Tông - HCM 2026) Cho hàm số $y = f(x)$ liên tục trên $\mathbb{R}$ và có bảng xét dấu đạo hàm như hình vẽ bên dưới. Hàm số đã cho có bao nhiêu điểm cực đại?\image{0.6}{}{lt_29.png}}{3}{1}{Đạo hàm đổi dấu 1 lần từ dương sang âm và $y = f(x)$ liên tục trên $\mathbb{R}$ nên hàm số có 1 điểm cực đại.}
	\option{$3$.}
	\option{$2$.}
	\option{$0$.}
	\option{$1$.}
\end{mcq}

\begin{mcq}{(THPT Lê Thánh Tông - HCM 2026) Cho hàm số $y = f(x)$ liên tục trên $[-1;+\infty)$ và có đồ thị như hình vẽ. Tìm giá trị lớn nhất của hàm số $y = f(x)$ trên $[1;4]$.\image{0.6}{}{lt_30.png}}{3}{1}{Giá trị lớn nhất của hàm số $y = f(x)$ trên $[1;4]$ là $y = 3$ đạt khi $x = 1$ và $x = 3$.}
	\option{$0$.}
	\option{$1$.}
	\option{$4$.}
	\option{$3$.}
\end{mcq}

\begin{mcq}{(THPT Lê Thánh Tông - HCM 2026) Biết rằng đồ thị hàm số $y = \frac{ax+1}{bx-2}$ có tiệm cận đứng là $x=2$ và tiệm cận ngang là $y=3$. Hiệu $a-2b$ có giá trị là}{2}{1}{$y = \frac{ax+1}{bx-2}$. Ta có tiệm cận đứng là $x = \frac{2}{b} = 2 \Rightarrow b = 1$.\\Tiệm cận ngang là $y = \frac{a}{b} = a = 3 \Rightarrow a = 3$.\\Suy ra $a - 2b = 3 - 2 = 1$.}
	\option{$4$.}
	\option{$0$.}
	\option{$1$.}
	\option{$5$.}
\end{mcq}

\end{quiz}

\end{document}
